\section{Expectation as a weighted average}

A probability mass function tells us how probability is distributed over the
possible values of a random variable. The next question is how to summarize that
whole distribution by a single representative number. The most important such
summary is expectation.

Expectation should not be introduced as a formula to memorize. It is a weighted
average. Values that occur with larger probability receive larger weight.

\subsection*{Averages with unequal weights}

Suppose a random variable takes values
\[
x_1,x_2,\ldots,x_n
\]
with probabilities
\[
p_1,p_2,\ldots,p_n,
\]
where
\[
p_i\geq 0,\qquad p_1+p_2+\cdots+p_n=1.
\]
If all values were equally likely, their average would be
\[
\frac{x_1+x_2+\cdots+x_n}{n}.
\]
But if some values are more likely than others, the average should count those
values more strongly. This leads to the weighted average
\[
x_1p_1+x_2p_2+\cdots+x_np_n.
\]

\begin{definitionbox}
For a discrete random variable \(X\) with probability mass function \(p_X\), the
\textbf{expectation} of \(X\) is
\[
\E[X]=\sum_{x\in\operatorname{supp}(X)}x\,p_X(x),
\]
provided the sum is well-defined.
\end{definitionbox}

For finite support, the sum is always finite. For countably infinite support,
one must check convergence; at this stage we mainly work with examples where the
sum is meaningful.

\subsection*{Expectation from the sample space}

If the original sample space is finite, expectation can also be computed by
summing over elementary outcomes:
\[
\E[X]=\sum_{\omega\in\Omega}X(\omega)P(\{\omega\}).
\]
This formula averages the numerical value \(X(\omega)\) over all elementary
outcomes, weighted by their probabilities.

The formula using the mass function is obtained by grouping together all
outcomes that produce the same value. If \(X=x\), then all outcomes in the event
\(\{X=x\}\) contribute the same value \(x\). The total probability of that event
is \(P(X=x)=p_X(x)\). Thus the grouped contribution is \(x p_X(x)\).

\begin{remarkbox}
The expectation can be computed either by averaging over elementary outcomes or
by averaging over values of the random variable. The second method is usually
more compact because it uses the distribution of \(X\).
\end{remarkbox}

\subsection*{Example: one fair die}

Let \(X\) be the result of one fair die roll. Then
\[
P(X=k)=\frac16,\qquad k=1,2,3,4,5,6.
\]
Therefore
\[
\E[X]=1\cdot\frac16+2\cdot\frac16+3\cdot\frac16+4\cdot\frac16+5\cdot\frac16+6\cdot\frac16.
\]
Hence
\[
\E[X]=\frac{1+2+3+4+5+6}{6}=\frac{21}{6}=3.5.
\]
The expectation is not necessarily a value that the random variable can take.
A die never shows \(3.5\). The expectation is a center of balance of the
probability distribution.

\subsection*{Example: payoff of a game}

Suppose a game has payoff random variable \(X\) defined by
\[
X=10 \quad\text{with probability }\frac16,
\]
\[
X=2 \quad\text{with probability }\frac26,
\]
and
\[
X=-3 \quad\text{with probability }\frac36.
\]
Then
\[
\E[X]=10\cdot\frac16+2\cdot\frac26+(-3)\cdot\frac36.
\]
Thus
\[
\E[X]=\frac{10}{6}+\frac{4}{6}-\frac{9}{6}=\frac56.
\]
The expected payoff is positive, even though some outcomes produce losses.
Interpreted over many repetitions, \(5/6\) is the average payoff per play that
the model predicts in the long run.

\subsection*{Expectation of an indicator}

Let \(I_A\) be the indicator of an event \(A\). Then
\[
I_A=1 \quad\text{with probability }P(A),
\]
and
\[
I_A=0 \quad\text{with probability }1-P(A).
\]
Therefore
\[
\E[I_A]=1\cdot P(A)+0\cdot(1-P(A))=P(A).
\]

\begin{theorembox}
For every event \(A\),
\[
\E[I_A]=P(A).
\]
\end{theorembox}

This identity is simple but fundamental. It says that probabilities can be read
as expectations of indicator variables.

\subsection*{Linearity of expectation}

Expectation has an important algebraic property: it is linear. If \(X\) and
\(Y\) are random variables and \(a,b\) are constants, then
\[
\E[aX+bY]=a\E[X]+b\E[Y].
\]
In finite models, this follows directly from the sample-space formula:
\[
\E[aX+bY]
=
\sum_{\omega\in\Omega}(aX(\omega)+bY(\omega))P(\{\omega\}).
\]
Distributing the sum gives
\[
\E[aX+bY]
=
a\sum_{\omega\in\Omega}X(\omega)P(\{\omega\})
+
b\sum_{\omega\in\Omega}Y(\omega)P(\{\omega\}).
\]
Hence
\[
\E[aX+bY]=a\E[X]+b\E[Y].
\]

The remarkable point is that linearity does not require independence. It is a
property of averaging.

\subsection*{Counting successes by indicators}

Toss a coin \(n\) times. Let
\[
X=\text{the number of heads}.
\]
Define indicator variables
\[
I_k=
\begin{cases}
1, & \text{the }k\text{-th toss is heads},\\
0, & \text{the }k\text{-th toss is tails}.
\end{cases}
\]
Then
\[
X=I_1+I_2+\cdots+I_n.
\]
If each toss has probability \(p\) of heads, then
\[
\E[I_k]=p
\]
for every \(k\). By linearity,
\[
\E[X]=\E[I_1]+\cdots+\E[I_n]=np.
\]

This derivation is often clearer than summing the binomial mass function. It
shows that the expected number of successes is the sum of the expected
contributions from each trial.

\subsection*{What expectation measures}

Expectation is not the most likely value. It is not necessarily a value that can
occur. It is a weighted average, or center of mass, of the distribution. In
applications, it often represents a long-run average under repeated independent
repetitions of the same model.

\begin{summarybox}
For a discrete random variable, expectation is computed by
\[
\E[X]=\sum_x xP(X=x).
\]
It summarizes the distribution by averaging values according to their
probability weights. It is best understood as a consequence of the probability
law of \(X\), not as a separate rule detached from the model.
\end{summarybox}
